\documentclass[11pt]{article}

\usepackage[margin=1in]{geometry}
\usepackage[T1]{fontenc}
\usepackage{lmodern}
\usepackage{hyperref}

\title{Annotated Bibliography for SymbolicLongMemorySequences.jl}
\author{Matthew Roughan and collaborators}
\date{\today}

\begin{document}
\maketitle

\section{Purpose}

This document records the literature foundation for SymbolicLongMemorySequences.jl, a Julia package for
synthesizing long-range dependent symbolic sequences. The accompanying BibTeX
database is \path{references/s5_references.bib}. Downloaded PDFs are stored
outside this public repository in the private sibling repository described in
\path{DOWNLOAD.mn}; this document and \path{references/README.md} retain
human-readable provenance and interpretation.

The BibTeX file has also been checked with PaperFetch.jl in online
check-only mode. Clear DOI/title/author conflicts were corrected in the
database, while low-confidence provider failures for books, proceedings, and
known local alternatives remain recorded as audit items rather than overwritten.

\section{Long-Range Dependence and Self-Similarity}

The general mathematical background comes from Samorodnitsky's survey
\cite{samorodnitsky2006_long}, Pipiras and Taqqu's monograph
\cite{pipiras2017_long}, and Beran's statistical treatment
\cite{beran1994_statistics}. These references distinguish covariance-based
long memory, spectral poles, heavy tails, self-similarity, and finite-sample
estimation. SymbolicLongMemorySequences.jl uses those distinctions to avoid presenting a finite
simulation as proof of asymptotic LRD.

The network-traffic literature, especially Leland et al.
\cite{leland1994_on} and Willinger, Paxson, and Taqqu
\cite{willinger1998_self}, motivates the practical need for synthetic data that
preserves large-scale burstiness. Taqqu, Teverovsky, and Willinger
\cite{taqqu1995_estimators}, Witt and Malamud \cite{witt2013_quantification},
and Malamud and Turcotte \cite{malamud1999_self} are used as cautions about
estimator dependence, finite scale ranges, trends, and diagnostic choices.

\section{Numerical LRD Synthesis}

Paxson's fast spectral synthesis \cite{paxson1997_fast} is the direct ancestor
of PB1 and the default numerical source for PB2 and PB3. Dieker
\cite{dieker2004_simulation}, Coeurjolly \cite{coeurjolly2000_simulation},
Wood and Chan \cite{wood1994_simulation}, Dietrich and Newsam
\cite{dietrich1997_fast}, and Craigmile \cite{craigmile2000_simulating}
document exact and approximate Gaussian alternatives. These references support
the current SymbolicLongMemorySequences.jl design choice: keep latent numerical generation explicit, so
future exact covariance backends can be added without changing the symbolic
transformation API.

Sorbye, Myrvoll-Nilsen, and Rue \cite{sorbye2017_an} show another scalable path
through finite mixtures of short-memory components. Roughan, Veitch, and Abry
\cite{roughan2000_real} provide the wavelet and multiscale background behind
PB3, which SymbolicLongMemorySequences.jl modifies into a latent-driver plus Markov-regime symbolic
generator.

\section{Point Processes and Heavy-Tailed Models}

Heavy-tailed and fractal point-process generators are represented by Garrett and
Willinger \cite{garrett1994_analysis}, Lowen and Teich
\cite{lowen1995_estimation}, Thurner et al. \cite{thurner1997_analysis}, and
Ryu and Lowen \cite{ryu1998_point}. SymbolicLongMemorySequences.jl adapts this line in MB2 through
heavy-tailed regime sojourns, and in MB3 through competing symbol-specific
renewal streams. Roughan, Yates, and Veitch \cite{roughan1999_the} motivate the
validation policy for these methods: finite simulations can miss the intended
scale range, so plots must mark interpretation limits.

\section{Symbolic and Categorical Applications}

Voss \cite{voss1992_evolution} is central because it analyzes DNA through one
binary indicator sequence per base, avoiding arbitrary scalar encodings of
symbols. SymbolicLongMemorySequences.jl's centered one-hot validation diagnostics follow that
methodological lesson. DNA references including Peng et al.
\cite{peng1992_long}, Li and Kaneko \cite{li1992_long}, Li, Marr, and Kaneko
\cite{li1994_understanding}, Buldyrev et al. \cite{buldyrev1998_analysis}, and
Buldyrev \cite{buldyrev2006_power} motivate symbolic biological use cases and
copy/mutation generator families.

Text references include Schenkel, Zhang, and Zhang \cite{schenkel1993_long},
Ebeling and Poschel \cite{ebeling1994_entropy}, Montemurro and Pury
\cite{montemurro2002_long}, Alvarez-Lacalle et al.
\cite{alvarez2006_hierarchical}, Altmann et al. \cite{altmann2009_beyond},
Altmann, Cristadoro, and Degli Esposti \cite{altmann2012_on},
Tanaka-Ishii and Bunde \cite{tanaka2016_long}, Debowski
\cite{debowski2015_hilberg}, and Wieczynski and Debowski
\cite{wieczynski2025_long}. These papers suggest mechanisms such as hierarchy,
rare-word clustering, topic persistence, and semantic similarity. SymbolicLongMemorySequences.jl does
not attempt to model text directly; it provides reproducible synthetic
sequences for studying estimators and sequence models under controlled symbolic
LRD.

\section{Model Details Used by SymbolicLongMemorySequences.jl}

PB1 combines Paxson-style fGn \cite{paxson1997_fast} with rank quantization.
PB2 combines fGn streams with an argmax transform inspired by latent Gaussian
categorical models \cite{gal2015_latent}. PB3 uses a latent spectral or Haar
driver to select Markov regimes, preserving direct local transition control.
PB4 uses an intermittent-map source motivated by Provata and Beck
\cite{provata2012_coupled}.

MB1a and MB1b are based on LAMP \cite{kumar2017_linear}; SymbolicLongMemorySequences.jl adds explicit
finite-history and dyadic-bucket variants. MB1c follows additive Markov-chain
memory-function work \cite{melnyk2006_memory,mayzelis2006_additive,melnik2014_entropy}.
MB2 and MB3 adapt heavy-tailed traffic and point-process models
\cite{garrett1994_analysis,ryu1998_point,lowen1995_estimation}. MB4 is
motivated by Hawkes-process word-occurrence modeling \cite{ogura2022_modeling}.
MB5 is motivated by symbolic DNA growth models
\cite{koroteev2015_duplication,salgado2012_exact}.

\bibliographystyle{IEEEtran}
\bibliography{references/s5_references}

\end{document}
